\section{Preguntas}

\addcontentsline{toc}{subsection}{Pregunta 1}

\subsection*{Pregunta 1}
\textbf{¿Que otros SMBD exiten actualmente en el mercado?:}

Un Sistema Manejador de Base de Datos (SMBD) es software orientado a comunicar a los usuarios y aplicaciones con las bases de datos, permitiendo así las modificaciones y consultas de la información almacenada, a su vez manteniendo la seguridad e integridad de ésta de forma estructurada y eficiente. 

Algunos otros SMBDs son, junto con algunas diferencias con respecto al instalado en esta práctica: \textit{Oracle Database}, \textit{Microsoft SQL Server}, \textit{MongoDB}, \textit{Redis} y \textit{SQLite}. Sus principales diferencias radican en el costo de licenciamiento, la velocidad provista, los principales propósitos que tienen, el público objetivo y el soporte a ciertos tipos. Por ejemplo, algunas de las muchas diferencias son:

\textbf{MySQL}
\begin{itemize}
        \item Cuenta con una implementación e interfaz de usuario más amigable.
    \item No admite tipos de datos más avanzados (geométricos, enumerados, matrices, rangos, XML, direcciones de red, hstore), sino que se limita a tipos numéricos, de caracteres, de fecha y hora, espaciales y JSON.
\end{itemize}
\textbf{MongoDB}
\begin{itemize}
    \item MongoDB Utiliza documentos JSON serializados con pares clave-valor, mientras que PostgreSQL utiliza una fila (tupla) como unidad básica de almacenamiento.
    \item MongoDB tiene un menor nivel de madurez comunitaria; pese a esto, ofrece controladores para múltiples lenguajes de programación.
    \item No soporta una variedad tan extensa de comandos SQL y tipos de datos (como arrays, json y hstore) como PostgreSQL.
\end{itemize}
\textbf{Microsoft SQL Server}
\begin{itemize}
    \item Es un producto comercial desarrollado por Microsoft que implica altos costos por licencias, soporte y funciones avanzadas.
    \item Ofrece mayor velocidad por sus capacidades de bases de datos en memoria y alto rendimiento.
\end{itemize}
\newpage
